% GENERATED FILE. Edit canonical lessons, then run npm run generate:books.
% canonical-source-hash: fnv1a64:f99ef8d98616f674
% canonical-lessons: PT-C421-escala, PT-C421-embarque, PT-C421-vaivem, PT-C421-pista, PT-C421-fuso

\chapter{No aeroporto}
\label{ch:pt-no-aeroporto}
\hlchaptermodality{421}

\begin{chapteropening}
\textbf{By the end of this chapter:} \emph{I can find my way through an airport: the stopover, boarding, the shuttle, the runway and the time zone.}

Complete the last lesson of chapter 421: I can find my way through an airport: the stopover, boarding, the shuttle, the runway and the time zone.
\end{chapteropening}

\section[a escala]{a escala — the stopover; the scale}

\label{lesson:PT-C421-escala}



\begin{quote}
Before the new one: say the Portuguese for seasickness, then the Portuguese for to feel sick.
\end{quote}

\subsection*{You'll want to know: a escala}
\textbf{a escala} — \textquotedblleft{}the stopover; the scale\textquotedblright{}.

\begin{quote}
O voo tem uma escala de três horas.
\end{quote}

\begin{quote}
\emph{The flight has a three-hour stopover.}
\end{quote}

\emph{Um voo sem escalas} — a direct flight. On a map, \emph{a escala} is the scale.

\begin{cousinweb}
\emph{escala} is a learned borrowing from Latin \emph{scāla} 'ladder, flight of steps' (EtymDB), from \emph{scandere} 'to climb'. A stopover is a step on the way.
\end{cousinweb}

\begin{grammarlens}[title={What you've built}]
One more word for telling, explaining and describing at B1.
\end{grammarlens}

\subsection*{Your turn}
\begin{itemize}
\raggedright
  \item \emph{Say it:} \emph{a escala}
  \item \emph{Say it:} \emph{a escala}, once more
  \item \emph{Say it:} a sentence of your own with \emph{a escala}
  \item \emph{Recall:} say the Portuguese for seasickness, then the Portuguese for to feel sick, then say \emph{a escala} again
\end{itemize}

\begin{tcolorbox}[breakable,colback=teal!4,colframe=teal!35!black,title={Before you move on}]
What does a escala mean? (\textquotedblleft{}The stopover; the scale\textquotedblright{}.) Say it once more.
\end{tcolorbox}
\section[o embarque]{o embarque — the boarding}

\label{lesson:PT-C421-embarque}



\begin{quote}
Before the new one: say the Portuguese for to feel sick, then the Portuguese for the stopover.
\end{quote}

\subsection*{You'll want to know: o embarque}
\textbf{o embarque} — \textquotedblleft{}the boarding\textquotedblright{}.

\begin{quote}
O embarque começa meia hora antes do voo.
\end{quote}

\begin{quote}
\emph{Boarding starts half an hour before the flight.}
\end{quote}

\emph{A porta de embarque} — the gate.

\begin{cousinweb}
\emph{embarque} is built on \emph{embarcar} 'to board', on \emph{barca} 'boat', inherited from Late Latin \emph{barca} (Lewis \& Short).
\end{cousinweb}

\begin{grammarlens}[title={What you've built}]
One more word for telling, explaining and describing at B1.
\end{grammarlens}

\subsection*{Your turn}
\begin{itemize}
\raggedright
  \item \emph{Say it:} \emph{o embarque}
  \item \emph{Say it:} \emph{o embarque}, once more
  \item \emph{Say it:} a sentence of your own with \emph{o embarque}
  \item \emph{Recall:} say the Portuguese for to feel sick, then the Portuguese for the stopover, then say \emph{o embarque} again
\end{itemize}

\begin{tcolorbox}[breakable,colback=teal!4,colframe=teal!35!black,title={Before you move on}]
What does o embarque mean? (\textquotedblleft{}The boarding\textquotedblright{}.) Say it once more.
\end{tcolorbox}
\section[o vaivém]{o vaivém — the shuttle (bus); the coming and going}

\label{lesson:PT-C421-vaivem}



\begin{quote}
Before the new one: say the Portuguese for the stopover, then the Portuguese for the boarding.
\end{quote}

\subsection*{You'll want to know: o vaivém}
\textbf{o vaivém} — \textquotedblleft{}the shuttle (bus); the coming and going\textquotedblright{}.

\begin{quote}
Um vaivém leva os passageiros até ao avião.
\end{quote}

\begin{quote}
\emph{A shuttle bus takes the passengers to the plane.}
\end{quote}

Literally 'goes-comes'. \emph{O vaivém} of a street is its bustle.

\begin{cousinweb}
\emph{vaivém} is built from \emph{vai} 'goes' and \emph{vem} 'comes' (EtymDB): something that keeps going back and forth.
\end{cousinweb}

\begin{grammarlens}[title={What you've built}]
One more word for telling, explaining and describing at B1.
\end{grammarlens}

\subsection*{Your turn}
\begin{itemize}
\raggedright
  \item \emph{Say it:} \emph{o vaivém}
  \item \emph{Say it:} \emph{o vaivém}, once more
  \item \emph{Say it:} a sentence of your own with \emph{o vaivém}
  \item \emph{Recall:} say the Portuguese for the stopover, then the Portuguese for the boarding, then say \emph{o vaivém} again
\end{itemize}

\begin{tcolorbox}[breakable,colback=teal!4,colframe=teal!35!black,title={Before you move on}]
What does o vaivém mean? (\textquotedblleft{}The shuttle (bus); the coming and going\textquotedblright{}.) Say it once more.
\end{tcolorbox}
\section[a pista]{a pista — the runway; the track; the clue}

\label{lesson:PT-C421-pista}



\begin{quote}
Before the new one: say the Portuguese for the boarding, then the Portuguese for the shuttle.
\end{quote}

\subsection*{You'll want to know: a pista}
\textbf{a pista} — \textquotedblleft{}the runway; the track; the clue\textquotedblright{}.

\begin{quote}
O avião está parado na pista.
\end{quote}

\begin{quote}
\emph{The plane is standing on the runway.}
\end{quote}

\emph{A pista de dança} — the dance floor. \emph{Uma pista} in a mystery is a clue.

\begin{cousinweb}
\emph{pista} was borrowed from Italian \emph{pista} 'beaten track', which is usually linked with Latin \emph{pinsere} 'to pound' (Lewis \& Short), but that is not settled.
\end{cousinweb}

\begin{grammarlens}[title={What you've built}]
One more word for telling, explaining and describing at B1.
\end{grammarlens}

\subsection*{Your turn}
\begin{itemize}
\raggedright
  \item \emph{Say it:} \emph{a pista}
  \item \emph{Say it:} \emph{a pista}, once more
  \item \emph{Say it:} a sentence of your own with \emph{a pista}
  \item \emph{Recall:} say the Portuguese for the boarding, then the Portuguese for the shuttle, then say \emph{a pista} again
\end{itemize}

\begin{tcolorbox}[breakable,colback=teal!4,colframe=teal!35!black,title={Before you move on}]
What does a pista mean? (\textquotedblleft{}The runway; the track; the clue\textquotedblright{}.) Say it once more.
\end{tcolorbox}
\section[o fuso]{o fuso — the time zone (o fuso horário); the spindle}

\label{lesson:PT-C421-fuso}



\begin{quote}
Before the new one: say the Portuguese for the shuttle, then the Portuguese for the runway.
\end{quote}

\subsection*{You'll want to know: o fuso}
\textbf{o fuso} — \textquotedblleft{}the time zone (o fuso horário); the spindle\textquotedblright{}.

\begin{quote}
Mudamos de fuso horário e perdemos uma hora.
\end{quote}

\begin{quote}
\emph{We change time zones and lose an hour.}
\end{quote}

\emph{O fuso horário} — the time zone. A spindle for spinning wool is also \emph{um fuso}.

\begin{cousinweb}
\emph{fuso} is inherited from Latin \emph{fūsus} 'spindle' (EtymDB): the time zones on a globe have the shape of a spindle.
\end{cousinweb}

\begin{grammarlens}[title={What you've built}]
One more word for telling, explaining and describing at B1.
\end{grammarlens}

\subsection*{Your turn}
\begin{itemize}
\raggedright
  \item \emph{Say it:} \emph{o fuso}
  \item \emph{Say it:} \emph{o fuso}, once more
  \item \emph{Say it:} a sentence of your own with \emph{o fuso}
  \item \emph{Recall:} say the Portuguese for the shuttle, then the Portuguese for the runway, then say \emph{o fuso} again
\end{itemize}

\begin{tcolorbox}[breakable,colback=teal!4,colframe=teal!35!black,title={Before you move on}]
What does o fuso mean? (\textquotedblleft{}The time zone (o fuso horário); the spindle\textquotedblright{}.) Say it once more.
\end{tcolorbox}
